% DNA double helix geometry: rise, pitch, diameter, major/minor groove.
\begin{tikzpicture}[>=Stealth, scale=0.95]
  % Two intertwined sinusoids forming the helix backbone
  \def\helixA{0.85}    % amplitude (radius proxy)
  \def\period{3.4}     % cm per turn
  \def\xmax{8.0}
  \draw[engblue, thick, smooth, samples=120, domain=0:\xmax]
    plot (\x, {\helixA*sin(2*pi*\x/\period r)});
  \draw[engamber, thick, smooth, samples=120, domain=0:\xmax]
    plot (\x, {\helixA*sin(2*pi*\x/\period r + pi r)});
  % Base-pair rungs every 0.34 cm (= 0.34 nm / bp at 1 nm = 1 cm)
  \foreach \k in {0,1,...,23} {%
    \pgfmathsetmacro\xpos{0.17 + 0.34*\k}
    \pgfmathsetmacro\ya{\helixA*sin(2*pi*\xpos/\period r)}
    \pgfmathsetmacro\yb{\helixA*sin(2*pi*\xpos/\period r + pi r)}
    \draw[enggray, very thin] (\xpos,\ya) -- (\xpos,\yb);
  }
  % Annotations
  \draw[<->, thick] (0, -1.25) -- (3.4, -1.25)
    node[midway, fill=white, inner sep=1pt, font=\small, align=center]
    {pitch $3.4$ nm \\ ($10$ bp / turn)};
  \draw[<->, thick] (3.4, -1.25) -- (3.74, -1.25)
    node[below=2pt, font=\footnotesize, align=center] {rise \\ $0.34$ nm/bp};
  \draw[<->, thick] (-0.4, -\helixA) -- (-0.4, \helixA)
    node[midway, fill=white, inner sep=1pt, font=\small, rotate=90, align=center]
    {diameter $2$ nm};
  % Major / minor groove indication
  \node[engred, font=\footnotesize] at (1.7, 1.35) {major groove ($2.2$ nm wide)};
  \node[engred, font=\footnotesize] at (3.4, -0.05) {minor groove ($1.2$ nm wide)};
  % Strand polarity labels
  \node[engblue, anchor=west, font=\footnotesize] at (\xmax+0.05, 0.6) {$5' \to 3'$};
  \node[engamber, anchor=west, font=\footnotesize] at (\xmax+0.05, -0.6) {$3' \to 5'$};
\end{tikzpicture}
